\section{USAJMO 2021}
        \begin{enumerate}
        	\item (\textit{USAJMO 2021}) $\textbf{Note:}$ For any geometry problem whose statement begins with an asterisk $(*)$, the first page of the solution must be a large, in-scale, clearly labeled diagram. Failure to meet this requirement will result in an automatic 1-point deduction.


 Let $\mathbb{N}$ denote the set of positive integers. Find all functions $f : \mathbb{N} \rightarrow \mathbb{N}$ such that for positive integers $a$ and $b,$
 \[f(a^2 + b^2) = f(a)f(b) \text{ and } f(a^2) = f(a)^2.\]


 



\textbf{Solution}

I claim that the only function $f$ that satisfies the constraints outlined within the problem is the function $f(n) = 1$ for all positive integers $n$.

 We will proceed with strong induction. The base case is simple, as plugging $a=1$ into the second equation given within the problem gives $f(1)=f(1)^2$. Since $f(n)$ can only return a positive integer value, we have that $f(1)=1$.

 Now we proceed with the inductive step. If the next number $n$ is either a perfect square or can be represented as a sum of two perfect squares, then obviously $f(n)=1$, as it is either the product of two $f$-values that are both equal to $1$ from the inductive assumption or is the square of an $f$-value that is equal to $1$, again due to the inductive assumption. Otherwise, we can use the Sum of Two Squares Theorem, which tells us that $n$ has at least one prime in its prime factorization that is $3(\mod 4)$ and is raised to an odd power.

 Lemma 1: Given that $a^2+b^2=c^2$ and $f(b)=f(c)=1$, we then have $f(a)=1$.

 Proof: Note that the first condition in the problem tells us that $f(a^2+b^2)=f(a)f(b)$, or $f(c^2)=f(a)f(b)$. Using the second condition gives us $f(c)^2=f(a)f(b)$. Plugging in the values of $f(b)$ and $f(c)$ gives us that $f(a)=1$.

 Now we will attempt to repeatedly remove prime factors that are $3(mod 4)$ and taken to an odd power, and we will move from the largest prime down to the smallest prime that satisfies above conditions. The prime factors will be removed by constructing a Pythagorean Triple with the prime being the smallest leg in the form of $p,\frac{p^2-1}{2},\frac{p^2+1}{2}$(this will always work as $p\neq 2$(not 3 mod 4), and this method works via Lemma 1). We will then prove the ending numbers that we achieve via removing all the $(\text{mod } 4)$ primes(which I will refer to as "tips") are equal to 1 b/c they can be expressed as the sum of two squares or is a perfect square(Sum of Two Squares Theorem). For example, if we took the number $21$, we would first aim to remove the $7$ by splitting it into $24$ and $25$, so $21$ would become $72$ and $75$. $75$ is divisible by $3$ to an odd power, so we transform it into $100$ and $125$. These two don't have any divisors that are $3 (\text{mod } 4)$ raised to an odd power so we leave it alone, as sum of two squares will work on them or they are perfect squares. $72$ does have a prime factor that is $3(\text{mod } 4)$, but it is an even power so we leave it alone. In this case, the tips are $72$, $100$, and $125$. However, now we need to prove two key facts: using this Pythagorean Triple Method will never generate another $3(\text{mod } 4)$ prime that is bigger than the current one we are working on or create more primes or keep the same number of primes in the tips' prime factorizations., as otherwise it could cause an infinite cycle, and also we must prove when we use Sum of Two Squares/Perfect Square given in the second condition on the tips the square root of the square(s) used will never be greater than or equal to $n$.

 The first claim can be proved rather simply. Note that $\frac{p^2-1}{2}$ can have no prime factors greater than or equal to $p$, as it can be factored as $(p-1)(\frac{p+1}{2})$, which are both less than $p$ and are integers($p+1$ must be an integer due to $p$ must being odd(not equal to $2$)). For $\frac{p^2+1}{2}$, we can prove something a bit more general.

 Lemma 2: For any positive integer $n$, all prime factors of $n^2+1$ must be $1(\text{mod } 4)$.

 Proof: Note that if $n^2+1\equiv 0(\text{mod } p)$, then we also must have $(n^2+1)(n^2-1)=(n^4-1)\equiv 0(\text{mod } p)$, or $n^4 \equiv 1(\text{mod } p)$. Now we can apply Fermat's Little Theorem to obtain $n^{p-1}\equiv 1(\text{mod } p)$. Note that since $n$ and $n^2$ are obviously not $1(\text{mod } p)$, as $n^2\equiv -1 (\text{mod } p)$, we have that $p-1$ is a multiple of $4$, or $p \equiv 1(\text{mod } 4)$.

 We can simply apply Lemma 2 to $p^2+1$ as $2$ is obviously not $3(\text{mod } 4)$, and this means that a prime factor $3(\text{mod } 4)$ will never be generated from this term. This completes the first of our two claims.

 Now we proceed to the second of our two claims. Note that every time we use the method on $n$ based on prime $p$, we will multiply by around $\frac{p}{2}$, clearly less than $p$(if we multiply by $p$ we would get $p^2$ which is clearly greater than $\frac{p^2+1}{2}$). We will never have to use the same prime twice in our method, so at max in the end we will multiply $n$ by the product of a little less than all $3(mod 4)$ primes that divide it, which is less than $n$ itself for $n$ greater than or equal to $3$(the smallest 3 mod 4 prime), meaning that the largest number that we must use two squares on that is generated by out method is less than $n^2$. We need to prove that the two squares that sum to this are both less than $n$, which is quite trivial, as they are less than $\sqrt{n^2}$, which obviously means it must be less than $n$. This proves the second of our claims.

 This completes the second case of the inductive step, and therefore completes both the induction and the problem.~Solution by hyxue

 



\textbf{Solution 2 (Taken from Twitch Solves ISL)}

The Answer is $\boxed{f\equiv 1}$ which works but we want to prove that it's the only one. Claim: If $f(a)=f(b)=1$ and $a>b$, then $f(a^2-b^2)= f(2ab)= 1$. Proof: We can write $1=f(a)=f(b)= f(a^2+b^2)= \sqrt{f((a^2+b^2)^2)} = \sqrt{f((a^2-b^2)^2+ (2ab)^2)} = \sqrt{f(a^2-b^2)f(2ab)}$. We set it to $a=b=1$ and we get that $f(1)=f(2)=1$. Easily by Induction it shows f(n)=1. We take if $n=2k$ take $(u,v)=(k,1)$ which $2uv=n$. If $n=2k+1$, just take $(u,v)= (k+1,k)$ which $u^2-v^2=n$. Thus the only answer is $\boxed{f\equiv 1}$ and we are done.

 



\textbf{Solution 3 (clear Solution 2)}

The answer is $\boxed{f\equiv 1}$, which works.To show it is necessary, we first get $f(1)=f(1)^2$, so $f(1)=1$.Then, we get $f(2)=f(1^2 + 1^2)=f(1)^2 =1$.

 The following claim finishes the problem via induction:

 Claim: If $f(n)=f(n-1)=1$ for $n \ge 2$, we have $f(2n)=f(2n-1)=1$.

 Proof: Note that $f(n^2+1)=f(n)f(1)=1$, so 
 \[1=f(n^2+1)^2 = f((n^2+1)^2)=f(2n)f(n^2 -1)\] implies $f(2n)=1$.

 Note that $f(2n^2 -2n+1)=f(n)f(n-1)=1$, so 
 \[1=f(2n^2 - 2n+1)^2 = f((2n^2-2n+1)^2)=f(2n-1)f(2n^2 -2n)\] implies $f(2n-1)=1$. $\square$

 Thus, the only solution is $\boxed{f\equiv 1}$.

 


	\item (\textit{USAJMO 2021}) $\textbf{Note:}$ For any geometry problem whose statement begins with an asterisk $(*)$, the first page of the solution must be a large, in-scale, clearly labeled diagram. Failure to meet this requirement will result in an automatic 1-point deduction.


 Rectangles $BCC_1B_2,$ $CAA_1C_2,$ and $ABB_1A_2$ are erected outside an acute triangle $ABC.$ Suppose that
 \[\angle BC_1C+\angle  CA_1A+\angle AB_1B=180^{\circ}.\]Prove that lines $B_1C_2,$ $C_1A_2,$ and $A_1B_2$ are concurrent.


 



\textbf{Solution}


                    \begin{figure}[H]
                        \centering
                        \includegraphics[width=0.5\textwidth]{Solution/USAJMO_I_2021_2_Solution_Figure_1.png}
                    \end{figure}
                    Let $D$ be the second point of intersection of the circles $(AB_1B)$ and $(AA_1C).$ Then:
 \begin{align*} \angle ADB &= 180^\circ – \angle AB_1B,&\angle ADC &= 180^\circ – \angle AA_1C\\ \angle BDC &= 360^\circ – \angle ADB – \angle ADC\\ &= 360^\circ – (180^\circ – \angle AB_1B) – (180^\circ – \angle AA_1C)\\ &= \angle AB_1B + \angle AA_1C\\ \angle BDC + \angle BC_1C &= 180^\circ \end{align*}
Therefore, $BDCC_1B_2$ is cyclic with diameters $BC_1$ and $CB_2$, and thus $\angle CDB_2 = 90^\circ.$Similarly, $\angle CDA_1 = 90^\circ$, meaning points $A_1$, $D$, and $B_2$ are collinear.

 Similarly, the points $A_2, D, C_1$ and $C_2, D, B_1$ are collinear.

 (After USAMO 2021 Solution Notes – Evan Chen)

 \textbf{vladimir.shelomovskii@gmail.com, vvsss}

 


	\item (\textit{USAJMO 2021}) $\textbf{Note:}$ For any geometry problem whose statement begins with an asterisk $(*)$, the first page of the solution must be a large, in-scale, clearly labeled diagram. Failure to meet this requirement will result in an automatic 1-point deduction.


 An equilateral triangle $\Delta$ of side length $L>0$ is given. Suppose that $n$ equilateral triangles with side length 1 and with non-overlapping interiors are drawn inside $\Delta$, such that each unit equilateral triangle has sides parallel to $\Delta$, but with opposite orientation. (An example with $n=2$ is drawn below.)

                    \begin{figure}[H]
                        \centering
                        \includegraphics[width=0.5\textwidth]{Problem/USAJMO_I_2021_3_Problem_Figure_1.png}
                    \end{figure}
                    
Prove that
 \[n \leq \frac{2}{3} L^{2}.\]


 



\textbf{Solution}

I will use the word "center" to refer to the centroid of any equilateral triangle.

 Consider the center of each triangle we place and the surrounding area that is "denied" placement for any centers. With a little bit of testing and sliding the equilateral triangle along each side of the center triangle, it is easy to see that this region forms a hexagon of side length $1$ where any other center being in this hexagon would mean that the two equilateral triangles at those centers would overlap. 

 Let us define a region for each center as its "personal space", where its personal space intersecting with any other center's personal space results in the two equilateral triangles intersecting. Given the above fact, it is easy to see that this personal space region encapsulates a hexagon of sidelength $\frac{1}{2}$ centered at each center by simply scaling the denied region by a factor of $\frac{1}{2}$. The area of each personal space would simply be $6(\frac{\sqrt{3}}{4}(\frac{1}{2})^2)=\frac{3\sqrt{3}}{8}$. Since each personal space cannot intersect, the maximum amount of equilateral triangles would be bounded by $\frac{\text{Area of Big Triangle}}{\text{Area of individual personal space}}=\frac{L^2\frac{\sqrt{3}}{4}}{\frac{3\sqrt{3}}{8}}=\frac{2}{3}L^2$, which completes the problem. (Unrigorous solution that I hope someone else can improve on ~ hyxue)

 


	\item (\textit{USAJMO 2021}) $\textbf{Note:}$ For any geometry problem whose statement begins with an asterisk $(*)$, the first page of the solution must be a large, in-scale, clearly labeled diagram. Failure to meet this requirement will result in an automatic 1-point deduction.


 Carina has three pins, labeled $A, B$, and $C$, respectively, located at the origin of the coordinate plane. In a move, Carina may move a pin to an adjacent lattice point at distance $1$ away. What is the least number of moves that Carina can make in order for triangle $ABC$ to have area 2021? 


 (A lattice point is a point $(x, y)$ in the coordinate plane where $x$ and $y$ are both integers, not necessarily positive.) 


 



\textbf{Solution 1}

The answer is $128$, achievable by $A=(10,0), B=(0,-63), C=(-54,1)$. We now show the bound.

 We first do the following optimizations:

 -if you have a point goes both left and right, we may obviously delete both of these moves and decrease the number of moves by $2$.

 -if all of $A,B,C$ lie on one side of the plane, for example $y>0$, we shift them all down, decreasing the number of moves by $3$, until one of the points is on $y=0$ for the first time.

 Now we may assume that $A=(a,d)$, $B=(b,-e)$, $C=(-c,f)$ where $a,b,c,d,e,f \geq 0$. Note we may still shift all $A,B,C$ down by $1$ if $d,f>0$, decreasing the number of moves by $1$, until one of $d,f$ is on $y=0$ for the first time. So we may assume one of $(a,b)$ and $(d,f)$ is $0$, by symmetry. In particular, by shoelace the answer to 2021 JMO Problem 4 is the minimum of the answers to the following problems:

 
Case 1 (where $a=d=0$) if $wx-yz=4042$, find the minimum possible value of $w+x+y+z$.

 Case 2 (else) $wy+xy+xz=(w+x)(y+z)-wz=4042$, find the minimum possible value of $w+x+y+z$.

 
Note that $(m+n)^2=4mn+(m-n)^2$ so if $m+n$ is fixed then $mn$ is maximized exactly when $|m-n|$ is minimized. In particular, if $m+n \leq 127$ then $mn-op \leq mn \leq 63*64 = 4032 <4042$ as desired.

 


	\item (\textit{USAJMO 2021}) $\textbf{Note:}$ For any geometry problem whose statement begins with an asterisk $(*)$, the first page of the solution must be a large, in-scale, clearly labeled diagram. Failure to meet this requirement will result in an automatic 1-point deduction.


 A finite set $S$ of positive integers has the property that, for each $s \in S,$ and each positive integer divisor $d$ of $s$, there exists a unique element $t \in S$ satisfying $\text{gcd}(s, t) = d$. (The elements $s$ and $t$ could be equal.) Given this information, find all possible values for the number of elements of $S$.


 



\textbf{Solution}

The answer is $0$ (left out by problem author) or $2^n$ for any non-negative integer $n$.

 To construct $|S|=2^n$, take $2n$ distinct primes $p_1, q_1, p_2, q_2, \dots p_n, q_n$ and construct $2^n$ elements for $S$ by taking the product over all $n$ indices $i$ of either $p_i$ or $q_i$ for every $i$.

 Note that (from the problem statement) each element of $S$ has $|S|$ divisors.

 Now it suffices to show that there is no element $x \in S$ such that $x$ is divisible by more than one power of any prime.

 Suppose there does exist an element $x$ and a prime $p$ such that $p^2 \mid x$.

 This implies that there exists an element $c \in S$ that is divisible by $p$ but not $p^2$ by $\gcd (x,c)=p$.

 Exactly half of $c$'s divisors are not divisible by $p$, so it follows that exactly half of the elements in $S$ are not divisible by $p$.

 However, if $p^0 , p^1 , \dots p^k$ are the powers of $p$ dividing $x$ for some $k \ge 2$, only $\frac{1}{k+1}$ of the divisors of $x$ are not divisible by $p$.

 But this means that $\frac{1}{k+1}$ of the elements of $S$ are not divisible by $p$, contradiction.

 Therefore, $|S|$ must be some power of $2$.

 


	\item (\textit{USAJMO 2021}) $\textbf{Note:}$ For any geometry problem whose statement begins with an asterisk $(*)$, the first page of the solution must be a large, in-scale, clearly labeled diagram. Failure to meet this requirement will result in an automatic 1-point deduction.


 Let $n \geq 4$ be an integer. Find all positive real solutions to the following system of $2n$ equations: 
 \begin{align*} a_{1} &=\frac{1}{a_{2 n}}+\frac{1}{a_{2}}, & a_{2}&=a_{1}+a_{3}, \\ a_{3}&=\frac{1}{a_{2}}+\frac{1}{a_{4}}, & a_{4}&=a_{3}+a_{5}, \\ a_{5}&=\frac{1}{a_{4}}+\frac{1}{a_{6}}, & a_{6}&=a_{5}+a_{7}, \\ &\vdots \\ a_{2 n-1}&=\frac{1}{a_{2 n-2}}+\frac{1}{a_{2 n}}, & a_{2 n}&=a_{2 n-1}+a_{1} \end{align*}



 



\textbf{Solution}

For any odd $n$, we know $a_n = \frac{1}{a_{n-1}} + \frac{1}{a_{n+1}}$. Additionally, for any even $m$, $a_m = a_{m-1} + a_{m+1}$. Substituting gives $a_m = \frac{2}{a_m} + \frac{1}{a_{m-2}} + \frac{1}{a_{m+2}}$, other than $m=2$ and $m=2n$.

 Consider the sequence $b_k$, such that $a_{2k} = b_k$. Then, say out of the set $[b_1, b_2, \cdots, b_n]$ let the maximum value of an element in the set be $X$ and the minimum value be $x$. So, $x \le X$.

 Then, we also know $b_k = \frac{2}{b_k} + \frac{1}{b_{k-1}} + \frac{1}{b_{k+1}}$, plug in $x$ into $b_k$ to get $x = \frac{2}{x} + \frac{1}{b_{k-1}} + \frac{1}{b_{k+1}}$.

 We know $\frac{1}{b_k} \ge \frac{1}{X}$, so $x \ge \frac{2}{x} + \frac{2}{X}$. Similarly, plugging in $X$ and using an inequality gives $\frac{2}{x} + \frac{2}{X} \ge X$.

 This means that $x \ge \frac{2}{x} + \frac{2}{X} \ge X$, so $x \ge X$. We also showed earlier that $x \le X$. Thus, $x = X$. So, all $a_{2k}$ are constants and equal to each other. 

 Now, let $a_{2k-1}=w$ and $a_{2k}=z$, then $w = \frac{2}{z}$ and $z=2w$. Solving gives $w=1$ and $z=2$

 


\end{enumerate}